\begin{tikzpicture}[font=\small,>=stealth,line join=round]
\draw[draw=teal!45!black,fill=green!3,rounded corners=3pt] (0,0) rectangle (13.6,6.8);
\node at (6.8,6.25) {导函数关于 $x=1$ 对称，原函数的中心高度由 $f(1)=7$ 确定};
\node at (3.5,5.55) {$f^{\prime}(x)=3(x-1)^2+2$};
\node at (10.1,5.55) {$f(x)=(x-1)^3+2(x-1)+7$};
\begin{scope}[shift={(1.6,1.4)},x=1.15cm,y=0.4cm]
\draw[gray,->] (-0.2,0)--(3.6,0) node[right] {$x$};
\draw[gray,->] (0,-0.3)--(0,9.5);
\draw[blue!70!black,very thick,domain=-0.2:2.2,samples=100] plot(\x,{3*(\x-1)^2+2});
\draw[blue!65!black,dashed] (1,0)--(1,8.7);
\fill[blue!70!black] (0,5) circle(2.5pt);
\fill[blue!70!black] (2,5) circle(2.5pt);
\draw[dotted,gray] (0,5)--(2,5);
\node[above] at (0,5.7) {$5$};\node[above] at (2,5.7) {$5$};
\draw[gray] (0,0+.08)--(0,0-.08);\node[below] at (0,0-.15) {$0$};
\draw[gray] (1,0+.08)--(1,0-.08);\node[below] at (1,0-.15) {$1$};
\draw[gray] (2,0+.08)--(2,0-.08);\node[below] at (2,0-.15) {$2$};
\end{scope}
\begin{scope}[shift={(8.9,1.4)},x=1.15cm,y=0.32cm]
\draw[gray,->] (-0.35,0)--(3.4,0) node[right] {$x$};
\draw[gray,->] (0,-0.3)--(0,11.7);
\draw[teal!75!black,very thick,domain=-0.2:2.2,samples=100] plot(\x,{(\x-1)^3+2*(\x-1)+7});
\fill[teal!75!black] (0,4) circle(2.5pt);
\fill[blue!70!black] (1,7) circle(2.5pt);
\fill[teal!75!black] (2,10) circle(2.5pt);
\draw[dotted,gray] (0,4)--(2,10);
\node[left] at (-.12,4) {$4$};\node[above right] at (2,10) {$10$};
\node[below right] at (1.1,6.1) {$(1,7)$};
\draw[gray] (0,0+.08)--(0,0-.08);\node[below] at (0,0-.15) {$0$};
\draw[gray] (1,0+.08)--(1,0-.08);\node[below] at (1,0-.15) {$1$};
\draw[gray] (2,0+.08)--(2,0-.08);\node[below] at (2,0-.15) {$2$};
\end{scope}
\node at (6.8,.45) {$0,2$ 关于 $1$ 配对；$f(0)+f(2)=4+10=14$。};
\end{tikzpicture}
